\documentclass[11pt,a4paper]{article}
\usepackage{fontspec}
\usepackage{xeCJK}
\setmainfont{msyh.ttc}[Path=C:/Windows/Fonts/,FontIndex=0,BoldFont=msyhbd.ttc]
\setCJKmainfont{msyh.ttc}[Path=C:/Windows/Fonts/,FontIndex=0,BoldFont=msyhbd.ttc]
\newfontfamily\titlelatin{simhei.ttf}[Path=C:/Windows/Fonts/,AutoFakeBold=2]
\newCJKfontfamily\titlecjk{simhei.ttf}[Path=C:/Windows/Fonts/,AutoFakeBold=2]
\usepackage[margin=22mm]{geometry}
\usepackage{graphicx}
\usepackage{needspace}
\usepackage{longtable}
\usepackage{array}
\usepackage[unicode,hidelinks,breaklinks]{hyperref}
\usepackage{fancyhdr}
\pagestyle{fancy}\fancyhf{}
\fancyfoot[C]{SRP 固定引导与原生自适应\quad\thepage}
\renewcommand{\headrulewidth}{0pt}
\setlength{\parindent}{0pt}
\setlength{\parskip}{7pt}
\linespread{1.28}
\emergencystretch=3em
\raggedbottom
\newcommand{\reviewtitle}[1]{\begin{center}{\titlelatin\titlecjk\bfseries\fontsize{21}{28}\selectfont #1}\end{center}\vspace{8pt}}
\newcommand{\reviewsection}[1]{\par\Needspace{9\baselineskip}\vspace{10pt}\begin{center}{\titlelatin\titlecjk\bfseries\fontsize{14}{20}\selectfont #1}\end{center}}
\newcommand{\reviewsubsection}[1]{\par\Needspace{4\baselineskip}\vspace{4pt}\begin{center}{\titlelatin\titlecjk\bfseries\fontsize{12}{18}\selectfont #1}\end{center}}
\begin{document}

\reviewtitle{SRP项目基石 · 科学依据与近邻反证}
设计候选 SRP-\allowbreak{}FIXED-\allowbreak{}VS-\allowbreak{}ADAPTIVE-\allowbreak{}NATIVE-\allowbreak{}20261008｜九篇近邻定向全文核对｜不宣称系统查新或增量已成立。


\reviewsection{一、先承认已有工作，再确定本项目增量}
\reviewsubsection{1.1 本轮重构对文献判断的影响}
用户明确的对象为固定引导与数据驱动原生自适应整包比较；B允许受限调整实际时序后回到共同理想。因此最近工作必须覆盖固定/\allowbreak{}反馈/\allowbreak{}自适应比较、自然反馈与可用性，不能沿用上一轮只关注艺术与抽象表示的差异判断。


新增全文核对ViBreathe、CHI 2026手戴自适应系统、Gummidela等2021游戏/\allowbreak{}引导比较及Prpa等2018，与既有四文合计八篇。七篇PDF已本地保存，Gummidela通过出版方完整HTML方法与结果阅读；文件下载失败不被写成已保存。PDF页码为物理页；阅读范围集中于设计、条件、剂量、测量、结果与局限，不冒充真实双人筛选。


\reviewsubsection{1.2 已被近邻覆盖的宽泛主张}
固定与反馈方案比较、响应式与自适应节奏、逐步靠近理想节奏、自然形态生理反馈，以及参与感/\allowbreak{}负担/\allowbreak{}执行权衡均已有研究。上述词语、设备组合和画面替换不构成独立创新。候选增量只能从具体多步骤场景问题、共同理想与移动提示的区分、可靠结果及失败边界中获得；若更广查新已覆盖，收缩为复现和应用验证。


\reviewsection{二、最直接的控制与比较近邻}
\reviewsubsection{2.1 L1 ViBreathe：已存在的固定、响应式与自适应比较}
Yu等，2021，ViBreathe: Heart Rate Variability Enhanced Respiration Training for Workaday Stress Management via an Eyes-\allowbreak{}free Tangible Interface，IJHCI，DOI 10.1080/\allowbreak{}10447318.2021.1898827。PDF第8–11页控制/\allowbreak{}方法，第12–15页结果；24人被试内比较三种引导，每次3分钟，另4人一周办公部署。自适应根据HRV改变后续提示周期，讨论舒适、投入和工作间隙应用。


\href{https://research.chalmers.se/publication/523202/file/523202_Fulltext.pdf}{L1机构公开全文}


对SRP的约束：不能把“数据驱动、受限改变节奏、回归理想”作为首次。候选区别是视觉原生场景中的多步骤语义、场景状态约束，以及同一I与实际G两参照下的目标代价。触觉到视觉、人群和剂量的差别本身只是范围变化；必须证明具体设计问题和新证据。原文HRV控制参数不直接移植到本项目。


\reviewsubsection{2.2 L2 CHI 2026：自然界面与自适应已高度接近}
Ruotong Yu等，2026，Enabling Adaptive Cardio-\allowbreak{}Respiratory Biofeedback Training on Ubiquitous Hand-\allowbreak{}Worn Devices，DOI 10.1145/\allowbreak{}3772318.3790488。PDF第6–8页双因素适应与莲花/\allowbreak{}莲叶界面，第10–13页评价。呼吸跟随与HRV趋势共同影响控制；48人四组、每次5分钟，报告生理、负担与体验指标，跟随评价参照实际引导频率。


\href{https://pi.cs.tsinghua.edu.cn/wp-content/uploads/2026/04/3772318.3790488.pdf}{L2作者实验室全文}


对SRP的约束：自然形式、双信号、自适应与固定比较的组合均不能直接证明原创性。SRP候选研究共同理想结构中的停留/\allowbreak{}补吸等多步骤，以及“跟上G”与“达到I”的差异；必须用可复算事实证明这种差异对体验或设计选择有用，不能仅新增两个符号当方法贡献。本文实验的视觉、姿态与引导组合也提示SRP应坦率报告整包解释范围。


不复用该文短窗、阈值或HRV优化结论；其控制目标、人群、装置与动作结构与本项目不同。基于HRV的状态解释仍需本项目独立测量依据。


\reviewsubsection{2.3 L3 游戏与呼吸引导：体验收益不保证执行收益}
Gummidela、da Cunha Silva与Gutierrez-\allowbreak{}Osuna，2021，Evaluating the Role of Breathing Guidance on Game-\allowbreak{}Based Interventions for Relaxation Training，Frontiers in Digital Health，DOI 10.3389/\allowbreak{}fdgth.2021.760268。完整HTML的Experimental interventions、Experiment 1/\allowbreak{}2与Discussion。比较节拍/\allowbreak{}反馈及有无游戏，实验一30人含额外对照，实验二36人；报告使用、执行与体验存在不同方向的结果。


\href{https://www.frontiersin.org/journals/digital-health/articles/10.3389/fdgth.2021.760268/full}{L3出版方完整HTML}


对SRP的约束：固定引导对比反馈、引导融入环境、练习投入与执行权衡均有直接先例；“更愿意使用”不能代替理想目标达成。SRP两臂可回答整包采用问题，不能像因子研究那样分别估计游戏/\allowbreak{}原生与反馈主效应。若未来要拆机制，须另加条件并重算样本；不是事后回归消除视觉差异。


\reviewsection{三、场景、真实性与解释性的近邻}
\reviewsubsection{3.1 L4 Attending to Breath：原生反馈需要可理解的映射}
Prpa等，DIS 2018，Attending to Breath，DOI 10.1145/\allowbreak{}3196709.3196765。PDF第4–6页系统/\allowbreak{}程序、第7–10页体验。11人、6分钟探索，未理解映射的8人接受说明后再体验6分钟；微现象学访谈用于研究呼吸觉察与控制感。


\href{https://www.metacreation.net/projects/respire}{L4作者实验室项目与全文入口}


该研究支持将反馈含义和控制感作为形成性问题，不能证明本项目更优。两次暴露和教学存在顺序差异，不能将结果当成SRP两臂因果效应。原生场景中“自行悟出含义”不应成为设计要求；我们的中性教学、提示/\allowbreak{}实际区分用于保证方案可用。用户无法理解或反而追逐画面时，应修订载体而非把困惑解释成沉浸。


\reviewsubsection{3.2 L5 Chittaro等：真实性的比较与整包比较不同}
Chittaro等，IJHCS 2024，Virtual reality experiences for breathing and relaxation training: The effects of real vs. placebo biofeedback，DOI 10.1016/\allowbreak{}j.ijhcs.2024.103275。PDF第10–13页方法结果：35人组间比较真实与他人记录反馈，分析31人，使用分阶段VR和多信号，测量含状态量表及体验。


\href{https://air.uniud.it/bitstream/11390/1278804/1/1-s2.0-S1071581924000594-main.pdf}{L5机构全文}


SRP的A是明确的固定指导与预设场景，不能称作假反馈对照；B另有原生形式及适应时序变化。故SRP估计整包效应，不能借该文把差异归因于真实性。若未来要回答真实闭环本身的增量，需同形态固定/\allowbreak{}回放比较另作设计，并处理教学、可信度及呈现统计差别。


\reviewsubsection{3.3 L6 Xu与Cho：艺术化比较已有先导}
Xu与Cho，2023，Exploring Artistic Visualization of Physiological Signals for Mindfulness and Relaxation: A Pilot Study，arXiv:2310.14343。PDF第3–5页：10人艺术化与图形反馈被试内比较，PANAS、易用性与放松评价并非一致方向。预印本身份保留。


\href{https://arxiv.org/abs/2310.14343}{L6原始预印本}


支持将主观喜欢、情绪与正确理解分开；不支持天气、粒子或B的自适应效应。SRP核心不再是艺术与抽象形式比较，不从它的摘要提炼“原生更优”。


\reviewsubsection{3.4 L7 Respireal：自然控制与使用成本}
Respireal，2025，DOI 10.1145/\allowbreak{}3715668.3736338，DIS Companion/\allowbreak{}Work in Progress。PDF第2–5页：前期6人访谈、后期7人体验与访谈，呼吸用于自然交互和MR/\allowbreak{}VR转换。其体验探索并未构成SRP的固定/\allowbreak{}自适应整包检验。


\href{https://exertiongameslab.org/wp-content/uploads/2025/06/respireal_dis2025.pdf}{L7作者实验室全文}


它提示自然交互具有设计空间及注意代价；学校短时场景、两设备准备成本和多步控制仍待本项目验证。自然呈现不能自动等同低负担。


\reviewsubsection{3.5 L8 Heartscape：抽象圆形也有价值与误读}
Rook等，2026，Heartscape: Biofeedback Through Real-\allowbreak{}Time Abstract Heart Rate Visualization，IJHCI，DOI 10.1080/\allowbreak{}10447318.2026.2669040。PDF第14–15页方法、第23–25页讨论：27人约10分钟体验与定性访谈，涉及视觉解释与注意；作者说明量化前后效益验证有待后续。


\href{https://researchportal.northumbria.ac.uk/ws/portalfiles/portal/231316276/Heartscape_Biofeedback_Through_Real-Time_Abstract_Heart_Rate_Visualization.pdf}{L8正式公开全文}


固定圈不能因为“抽象”就被先验当作较差对照；B的复杂反馈也可能带来解释负担。其个体解释不能作为情绪真值或本项目量表效度。


\reviewsection{四、需求—原理—设计—证据矩阵}
\reviewsubsection{4.1 支持范围与取消条件}
E1 呼吸练习的总体动机。Fincham等2023荟萃的压力主结果涵盖12项RCT、785成人；Balban等2023研究为5分钟/\allowbreak{}天、一月的不同呼吸结构比较。它们支持进一步研究呼吸练习，不支持当前四天气、单次旅程或精确秒数。若目标情境没有可接受收益，则收缩产品效益表述，保留可靠的相对比较结果。


\href{https://www.nature.com/articles/s41598-022-27247-y}{呼吸练习荟萃原文}


\href{https://doi.org/10.1016/j.xcrm.2022.100895}{呼吸结构原始研究}


E2 自适应可能改善可跟随性，也可能降低要求。L1/\allowbreak{}L2提供近邻依据；本项目I/\allowbreak{}G/\allowbreak{}X分离是为避免评价循环。需要独立X、规则冻结及共同I的达成证据。若只提高G跟随率而I偏差增加，取消“训练目标更有效”的主张，报告支持与目标漂移的权衡。


E3 场景融入可能改变注意、控制感和情绪。L3/\allowbreak{}L4/\allowbreak{}L7/\allowbreak{}L8支撑可研究性，不给出原生必胜结论。需要中性教学、真实体验及错误案例。若注意主要停留在猜测系统或追逐奖励，收缩装饰及不必要反馈。


E4 实际数据只能支持可靠可辨的交互状态。HRV研究方法建议强调呼吸、记录窗口和伪迹影响；H10不能自动给出情绪状态。未来加入心率相关控制前，先证明控制变量的可识别性、时间尺度与额外作用；否则只保留辅助测量。


\href{https://www.frontiersin.org/journals/psychology/articles/10.3389/fpsyg.2017.00213/full}{Laborde等HRV方法建议}


E5 输入不可用时的诚实表现来自必要语义边界。B应停止基于无效X的适应并明确降级；A本来就是开环。需要故障与恢复材料。若断流时仍自动显示实际成功或累计达标，不能称可信闭环；质量图标存在不等于质量检测正确。


E6 四天气及Storm/\allowbreak{}Fade多步设计是应用和创作选择。现有3—3—3—3、4—6、5—5、2.5—1.5—6为工程候选；停留、补吸和时长必须经舒适性/\allowbreak{}可辨性验证。若结构对任务无必要价值或不可可靠识别，简化设计并重定研究范围，不因资产已有而强留。


\reviewsubsection{4.2 参数的确定方法}
适应范围、单次变化、观察窗、回归速率和缓冲预算：先用独立真实标注判断检测延迟/\allowbreak{}误差，再用获准的舒适性及任务理解预试确定可接受区域，最后冻结规则和参数。不是最大化正式样本得分；不同人可按同一预设校准程序获得参数。


主要量表、最小关注差、前后相关、残差与缺失：按产品用途和前置分布论证。理解题数不再用于主要N计算；旧PANAS量表基础可复用，但旧条件的效果和功效结果不能直接外推。功能容差应来自测量误差与有意义偏离，禁止让适应器与评分器共同放宽边界。


教学/\allowbreak{}核心时长及阶段人数是可行性候选，预试关注理解、真值、负担与实施条件，不用小样本显著性为正式效益放行。Eldridge等先导方法建议用于区分可行性与确认性目标，不提供本项目人数的普适阈值。


\href{https://www.bmj.com/content/355/bmj.i5239}{先导与可行性研究报告方法}


\reviewsection{五、贡献判断与证据缺口}
\reviewsubsection{5.1 从“有趣系统”到“新增知识”的条件}
原创设计工作包括团队的具体场景、多步骤控制语义及配对制品；学术新颖性需要证明相对于L1/\allowbreak{}L2/\allowbreak{}L3等最近工作解决了尚未解决的具体问题。C1侧重制品与方法的明确性，C2侧重可信权衡，C3侧重有限规则的支持与反例。增加天气或指标本身不增加贡献等级。


当前最值得检验的候选问题是：当原生闭环帮助用户跟上动态提示时，它是否仍帮助用户接近共同理想多步结构，还是仅通过任务放宽获得舒适与好看反馈？这种区分必须改变真实设计判断才有研究价值；I/\allowbreak{}G/\allowbreak{}X作为记录设计不是新的数学理论。


\reviewsubsection{5.2 继续查新与发表定位}
八篇为定向全文近邻，不是检索穷尽。W-\allowbreak{}01应继续追溯L1/\allowbreak{}L2/\allowbreak{}L3引用及近年多步骤、适应评价、情绪测量文献；Stairway to Heaven已补获出版社HTML全文，见6.9；其直接引用的自适应先例仍须逐篇核查。不能从某篇没报告某字段推断全领域空白。若多步骤与理想参照的增量也已覆盖，定位为有明确情境价值的复现/\allowbreak{}应用研究，重新判断期刊适配。


CCF B及以上为研究目标。比单纯展示更有力的证据是有效对照、独立真值、共同目标的效益/\allowbreak{}代价及可解释失败；本设计完成不保证等级、录用或实践奖项。比赛后续使用同一真实制品和价值证据，不能把动画联动冒充自适应已经实现。

% NORMALIZATION-20261008
\newpage\reviewsection{六、规范化的主张—证据对照}
\reviewsubsection{6.1 L1 ViBreathe}
定位：PDF8–11控制/\allowbreak{}方法，12–15结果。对象/\allowbreak{}剂量：24人被试内三提示；另4人办公部署；固定、响应、自适应；3分钟。

实际支持：自适应与固定比较已有；体验与执行不等同。本项目可用：支持有限调整及效益/\allowbreak{}代价问题。不能推导：不能证明视觉多步骤、新颖性或本项目参数。

\href{https://doi.org/10.1080/10447318.2021.1898827}{原始来源入口}

\reviewsubsection{6.2 L2 Yu CHI2026}
定位：PDF6–8控制/\allowbreak{}界面，10–13评价。对象/\allowbreak{}剂量：48人四组；5分钟；双信号适应与自然形态。

实际支持：自然形态与适应已有，跟随以实际提示为参照。本项目可用：要求分离G跟随与I达成。不能推导：不能把换载体/\allowbreak{}加I符号当新方法。

\href{https://doi.org/10.1145/3772318.3790488}{原始来源入口}

\reviewsubsection{6.3 L3 Gummidela2021}
定位：出版方HTML Experimental interventions /\allowbreak{} Experiment1,2 /\allowbreak{} Discussion。对象/\allowbreak{}剂量：实验一30、实验二36；节拍/\allowbreak{}反馈及有无游戏。

实际支持：使用、执行与体验并非同向。本项目可用：整包采用应分别报告收益与代价。不能推导：两臂不能拆游戏/\allowbreak{}原生/\allowbreak{}适应效应。

\href{https://doi.org/10.3389/fdgth.2021.760268}{原始来源入口}

\reviewsubsection{6.4 L4 Prpa2018}
定位：PDF4–6系统程序，7–10体验。对象/\allowbreak{}剂量：11人探索；6分钟；未理解者说明后再体验。

实际支持：映射理解和控制感需要解释。本项目可用：中性教学与误读材料。不能推导：非本项目两臂因果/\allowbreak{}情绪效果证据。

\href{https://doi.org/10.1145/3196709.3196765}{原始来源入口}

\reviewsubsection{6.5 L5 Chittaro2024}
定位：PDF10–13方法结果。对象/\allowbreak{}剂量：35分配、31分析；真实与他人记录反馈。

实际支持：真实性与多信号体验比较。本项目可用：提醒反馈真实性需独立核查。不能推导：A非假反馈；不能借其归因真实闭环机制。

\href{https://doi.org/10.1016/j.ijhcs.2024.103275}{原始来源入口}

\reviewsubsection{6.6 L6 XuCho2023}
定位：PDF3–5方法结果。对象/\allowbreak{}剂量：10人先导；艺术与图形反馈被试内。

实际支持：情绪、放松和易用性不同向。本项目可用：保留情绪/\allowbreak{}喜欢/\allowbreak{}理解独立维度。不能推导：预印本、小样本不证明B自适应收益。

\href{https://arxiv.org/abs/2310.14343}{原始来源入口}

\reviewsubsection{6.7 L7 Respireal2025}
定位：PDF2–5。对象/\allowbreak{}剂量：前期6访谈、后期7体验；自然呼吸交互；WIP。

实际支持：自然交互有注意与使用代价。本项目可用：校园场景仍需成本/\allowbreak{}需求验证。不能推导：非固定/\allowbreak{}适应整包确认性研究。

\href{https://doi.org/10.1145/3715668.3736338}{原始来源入口}

\reviewsubsection{6.8 L8 Heartscape2026}
定位：PDF14–15方法，23–25讨论。对象/\allowbreak{}剂量：27人体验访谈；约10分钟抽象心率可视化。

实际支持：抽象表示也有解释价值/\allowbreak{}误读。本项目可用：A不能预设为劣质，B复杂度需审查。不能推导：无量化前后收益证明，不提供呼吸控制真值。

\href{https://doi.org/10.1080/10447318.2026.2669040}{原始来源入口}

\reviewsubsection{6.9 L9 Stairway to Heaven: A Gamified VR Journey for Breath Awareness}
定位：2026-10-08经正常浏览器取得ACM开放HTML全文；PDF下载未成功。Nathan Miner等，CHI2024，DOI 10.1145/3613904.3641986。3.1—3.3为机制，4.1—4.3为方法，5.2—5.3为结果，6.2与7为边界。22人招募，21人问卷/访谈，18人呼吸分析；平均年龄30岁，非本项目校园样本。平均体验15.7分钟；单方案混合方法，无固定引导平行对照。

实际支持：校准呼吸深度门槛驱动HUD与旅程进展，用户自行决定节奏，系统不调节呼吸节奏（3.1、6.2）。问卷为ITQ/IPQ及改为七点评分的SUS，非PANAS；前3人后校准脚本修改，呼吸分析保留18人（4.1、4.3）。结果报告体验中的平均呼吸速率下降；同时发现校准负担、疲劳及场景/声音/呼吸的注意冲突。

本项目可用：自然反馈与游戏进展已有先例，撤回宽泛首次主张；本项目B对G的受限放慢与回归不同于该作品自行调速。将校准、门槛和疲劳纳入成本/失败解释。不能推导：单方案前中后变化不证明B优于A，未检验本项目共同I/G评分、短时负性情绪效益或独立迁移。机制差异只界定候选增量，不直接证明新颖性。该文2.2引用BreathCoach等自适应先例，不能据此宣称首次自适应。

\href{https://doi.org/10.1145/3613904.3641986}{原始来源入口}

\newpage\reviewsection{七、量表身份与设计选择的证据边界}
20题PANAS原版、黄丽等2003中国人群适用研究、邱林等2008的9+9修订及60题PANAS-\allowbreak{}X分别记录。黄丽文献在期刊站点的身份/摘要已核准。侧面取得CHEUNG 2018香港教育大学博士论文：方法印刷页41引用黄丽译本，中英附录印刷页127、133—134可核题序，但采用过去1—2周指令。此题本与2003原始译本的一致性及此刻适用仍待核准，不能给改写指令贴验证标签。邱林摘要确认18题修订；不套用NA10题计分。PhenX说明对应60题PANAS-\allowbreak{}X，不能外推本20题中文版本使用条件。

\href{https://doi.org/10.1037/0022-3514.54.6.1063}{原20题研究}；\href{https://doi.org/10.3321/j.issn:1000-6729.2003.01.018}{HYJ2003身份}；\href{https://cnki.istiz.org.cn/kcms/detail/detail.aspx?dbcode=CJFQ&dbname=CJFD2008&filename=YXNX200803009}{Qiu2008修订}；\href{https://www.phenxtoolkit.org/protocols/view/180501}{PANAS-X条件}

\href{https://www.lib.eduhk.hk/pure-data/pub/202201871.pdf}{CHEUNG 2018机构库题本侧证}；\href{https://zxws.cbpt.cnki.net/portal/journal/portal/client/paper/a793f765d726bf9a5f28aed99964ee07}{黄丽等2003期刊身份与摘要}。馆藏水印限制进一步复制与出版，本包只记题序/身份，不公开完整题文。

状态指令已有直接使用先例：\href{https://doi.org/10.1111/aphw.12133}{Hu等2018，PDF7}采用中文20题、此刻、五级评分及分量表求和；\href{https://doi.org/10.3724/SP.J.1041.2021.00244}{刘宇平等2021，PDF7}采用中文NA10题与此时此刻指令。两文全文已取得，但没有据此核准本项目逐字中文题本；人群/任务不同，不代证普通学习间隙的低NA分布与变化敏感性。后者属于前置真实验证。

\href{https://doi.org/10.1037/t03592-000}{APA原20题数据库记录}已直接核准身份和时间指令描述，许可字段尚未取得。\href{https://www.apa.org/pubs/databases/psyctests/faq}{官方FAQ}按具体记录区分非商业研究/教学与公开转载；馆藏水印不能替代具体量表的完整使用条件。M01保留题本/词义/来源对应缺口；M02分别核对施测、改编和公开用途。

\reviewsubsection{7.1 Storm}
3—3—3—3停留结构是工程候选；停留自愿，不强求。雨幕/\allowbreak{}风纹映射是创作选择，不证明雨天具有特定情绪效果。 用途/\allowbreak{}舒适/\allowbreak{}停留可辨不成立时，两组共同改目标或简化。

\reviewsubsection{7.2 Heat}
4—6吸呼比是工程候选，热景不用于判定学生情绪；较长呼气用途需结构文献与本场景核查。 无法解释额外作用则仅作实例或删减。

\reviewsubsection{7.3 Snow}
5—5是工程候选；低动态天气载体不能替代步骤真值，颜色不单独承担状态。 信息不可辨则简化载体。

\reviewsubsection{7.4 Fade}
2.5—1.5—6与补潮是工程/\allowbreak{}创作候选；不同于把其他剂量/\allowbreak{}结构原研究效果直接移植。 补吸无法独立判别则退出完整三步闭环，不追补。

\reviewsubsection{7.5 实时/累计}
实时反馈要求可靠X，累计只说明可靠G练习响应；文献支持可研究性，不保证收益。 误读/\allowbreak{}额外负担无法修复就取消累计/\allowbreak{}装饰。

\reviewsubsection{7.6 两设备}
respiBAN提供呼吸/\allowbreak{}运动原件，H10只辅助研究，不参与首版控制。佩戴相同保证研究测量条件，不证明部署必需。 没有额外研究用途就重新论证H10产品价值；本轮保留设备约束。


\end{document}
